% Source presentations for register row N-c-009.

\subsection{What the original \(\mathcal{ALC}\) completion calculus proves}
\label{sec:p09-presentation-n-c-009}

\subsubsection{Language, semantics, and exact decision problem}

Schmidt-Schau\ss--Smolka fix disjoint alphabets of concept symbols and role
symbols, reserve \(\top\) as a concept symbol, and generate their full
description language by
\[
 C,D ::= A
 \mid \forall R:C
 \mid \exists R:C
 \mid C\sqcap D
 \mid C\sqcup D
 \mid \neg C.
\]
The colon is the article's role-restriction notation; the dot in
Definition~\ref{def:p09-alc} expresses the same constructor.  Both union and
existential restriction occur in the printed grammar even though complement
makes them eliminable.  The name \(\mathcal{ALC}\) stands here for this
object-language of concept descriptions, not for a knowledge-base input format
\cite[\S2, pp.~5--6]{SchmidtSchaussSmolka1991}.

An interpretation is printed as
\(\mathcal I=(D^{\mathcal I},{}^{\mathcal I})\).  Concept symbols denote
subsets of \(D^{\mathcal I}\), role symbols denote subsets of
\(D^{\mathcal I}\times D^{\mathcal I}\), and
\[
\begin{aligned}
 \top^{\mathcal I}&=D^{\mathcal I},\\
 (\forall R:C)^{\mathcal I}
   &=\{a\in D^{\mathcal I}:
       \forall b\,((a,b)\in R^{\mathcal I}\Rightarrow
                         b\in C^{\mathcal I})\},\\
 (\exists R:C)^{\mathcal I}
   &=\{a\in D^{\mathcal I}:
       \exists b\,((a,b)\in R^{\mathcal I}\land
                         b\in C^{\mathcal I})\},\\
 (C\sqcap D)^{\mathcal I}&=C^{\mathcal I}\cap D^{\mathcal I},\\
 (C\sqcup D)^{\mathcal I}&=C^{\mathcal I}\cup D^{\mathcal I},\\
 (\neg C)^{\mathcal I}&=D^{\mathcal I}\setminus C^{\mathcal I}.
\end{aligned}
\]
The article says that \(D^{\mathcal I}\) is a set; it does not explicitly
print the now-standard condition \(D^{\mathcal I}\ne\varnothing\).  This
omission does not change its coherence problem, because a coherent concept has
a nonempty extension and therefore a nonempty witnessing domain.  It does
matter to the literal generality of the later model extraction for a constraint
system containing no individual variable.  Definition~\ref{def:p09-alc}
separately adopts the modern nonempty-domain convention and must not be read
back into the historical definition
\cite[\S2, p.~5]{SchmidtSchaussSmolka1991}.

For descriptions \(C,D\), the source writes
\[
 C\sim D
 \quad\Longleftrightarrow\quad
 C^{\mathcal I}=D^{\mathcal I}
 \text{ for every interpretation }\mathcal I.
\]
This is semantic equivalence between concept descriptions.  It is used to
distinguish the printed primitives from the source's derived concept
expressions:
\[
\begin{aligned}
 \top&\sim A\sqcup\neg A,\\
 \exists R:C&\sim\neg(\forall R:\neg C),\\
 C\sqcup D&\sim\neg(\neg C\sqcap\neg D),\\
 C-D&:=C\sqcap\neg D,\\
 C\mathbin{\to}D&:=\neg C\sqcup D.
\end{aligned}
\]
The last two lines introduce source abbreviations for concept difference and
concept implication, not new semantic categories.  Difference is used by the
Introduction's complement examples; implication lets a pointwise conditional
be written as a concept while remaining distinct from the metalinguistic
subsumption statement \(C\sqsubseteq D\)
\cite[pp.~2--6]{SchmidtSchaussSmolka1991}.

A description \(C\) is \emph{coherent} exactly when
\(C^{\mathcal I}\ne\varnothing\) in some interpretation.  It is
\emph{subsumed by} \(D\) exactly when
\[
 C^{\mathcal I}\subseteq D^{\mathcal I}
 \quad\text{for every interpretation }\mathcal I.
\]
Proposition~2.1 gives the two linear reductions
\[
\begin{aligned}
 C\text{ is subsumed by }D
 &\quad\Longleftrightarrow\quad
 C\sqcap\neg D\text{ is incoherent},\\
 C\text{ is incoherent}
 &\quad\Longleftrightarrow\quad
 C\text{ is subsumed by }\neg\top.
\end{aligned}
\]
The article's Introduction also uses the active phrase \(C\) \emph{subsumes}
\(D\) for the reverse containment \(D^{\mathcal I}\subseteq C^{\mathcal I}\).
Keeping the displayed argument order avoids that grammatical reversal
\cite[pp.~3, 5]{SchmidtSchaussSmolka1991}.

For one concept and no supplied TBox or ABox, historical coherence is exactly
the modern concept-satisfiability task used in
Theorem~\ref{thm:p09-alc-modal}.  The article's failure to print a nonempty
domain convention does not disturb this bridge: a coherence witness already
contains an element of \(C^{\mathcal I}\).

The source calls \(\neg A\), for a concept symbol \(A\), a \emph{simple
complement}; a description is \emph{simple} when every complement in it has
that form.  These terms name syntactic classes and are used by the constraint
translation.  Complex complements are pushed inward by
\[
\begin{aligned}
 \neg(\forall R:C)&\sim\exists R:\neg C,&
 \neg(\exists R:C)&\sim\forall R:\neg C,\\
 \neg(C\sqcap D)&\sim\neg C\sqcup\neg D,&
 \neg(C\sqcup D)&\sim\neg C\sqcap\neg D,\\
 \neg\neg C&\sim C.
\end{aligned}
\]
Proposition~2.2 states, without printing the induction or size calculation,
that this yields an equivalent simple description in linear time.  The source
also fixes three syntactic restrictions:
\[
\begin{array}{c|l}
 \mathcal{ALE}&\text{simple descriptions without union},\\
 \mathcal{ALU}&\text{simple descriptions whose existentials have form }
                       \exists R:\top,\\
 \mathcal{AL}&\mathcal{ALE}\cap\mathcal{ALU}.
\end{array}
\]
Proposition~2.3 reads propositional CNF as intersections of unions of concept
symbols and simple complements.  A satisfying valuation gives a one-element
interpretation, and a point in a coherent extension gives a satisfying
valuation.  Thus \(\mathcal{ALU}\)-coherence is NP-hard and
\(\mathcal{ALU}\)-subsumption is co-NP-hard even without roles
\cite[\S2, pp.~6--7]{SchmidtSchaussSmolka1991}.

These are all-interpretations problems for one concept description.  The
Introduction uses terminological equations and inclusions as motivation and
refers to earlier reductions under restrictions, but this article supplies no
general TBox, no ABox input, and no knowledge-base completion calculus.

\subsubsection{The six constraint forms and deterministic unfolding}

Section~3 uses individual variables \(x,y,z\) and concept variables \(X,Y,Z\),
with the two alphabets disjoint.  For a fixed interpretation \(\mathcal I\), an
assignment \(\alpha\) maps individual variables into \(D^{\mathcal I}\) and
concept variables to subsets of \(D^{\mathcal I}\).  The constraint language
has exactly the following six forms:
\[
 X\sqsubseteq C,\qquad
 X(\forall R)Y,\qquad
 X(\exists R)Y,\qquad
 X\sqsubseteq Y\sqcup Z,\qquad
 x:X,\qquad
 xRy,
\]
where \(C\) in the first form is a simple description.  Their satisfaction
clauses are
\[
\begin{aligned}
 \alpha\models_{\mathcal I}X\sqsubseteq C
 &\iff \alpha(X)\subseteq C^{\mathcal I},\\
 \alpha\models_{\mathcal I}X(\forall R)Y
 &\iff
   \forall a\in\alpha(X)\,
   \forall b\,((a,b)\in R^{\mathcal I}\Rightarrow b\in\alpha(Y)),\\
 \alpha\models_{\mathcal I}X(\exists R)Y
 &\iff
   \forall a\in\alpha(X)\,
   \exists b\,((a,b)\in R^{\mathcal I}\land b\in\alpha(Y)),\\
 \alpha\models_{\mathcal I}X\sqsubseteq Y\sqcup Z
 &\iff \alpha(X)\subseteq\alpha(Y)\cup\alpha(Z),\\
 \alpha\models_{\mathcal I}x:X
 &\iff \alpha(x)\in\alpha(X),\\
 \alpha\models_{\mathcal I}xRy
 &\iff (\alpha(x),\alpha(y))\in R^{\mathcal I}.
\end{aligned}
\]
A constraint system is a finite nonempty set \(S\) of constraints.  It is
satisfiable when one pair \((\mathcal I,\alpha)\) satisfies every member.
Consequently, for fresh \(x,X\), Proposition~3.1 states
\[
 C\text{ coherent}
 \quad\Longleftrightarrow\quad
 \{x:X,\ X\sqsubseteq C\}\text{ satisfiable}.
\]
The proposition is stated without proof
\cite[\S3, pp.~8--9]{SchmidtSchaussSmolka1991}.

A constraint system is \emph{simple} when its first-form constraints have only
\(X\sqsubseteq A\), with \(A\ne\top\), or \(X\sqsubseteq\neg A\).  The
following rewrites eliminate every other first-form right-hand side:
\[
\begin{array}{rcll}
 X\sqsubseteq\forall R:C
 &\rightsquigarrow&X(\forall R)Y,\ Y\sqsubseteq C,
 &Y\text{ fresh},\\
 X\sqsubseteq\exists R:C
 &\rightsquigarrow&X(\exists R)Y,\ Y\sqsubseteq C,
 &Y\text{ fresh},\\
 X\sqsubseteq C\sqcap D
 &\rightsquigarrow&X\sqsubseteq C,\ X\sqsubseteq D,\\
 X\sqsubseteq C\sqcup D
 &\rightsquigarrow&X\sqsubseteq Y\sqcup Z,\
                    Y\sqsubseteq C,\ Z\sqsubseteq D,
 &Y,Z\text{ fresh},\\
 X\sqsubseteq\top
 &\rightsquigarrow&\text{no constraint}.&
\end{array}
\]
Here \emph{simple description} names an object-language class determined by
the placement of complements, whereas \emph{simple constraint system} names a
normal form after translation.  Neither definition subsumes the other.
Fresh-variable choice changes only names.  These are deterministic
syntax-unfolding steps, not model-search branches.  Propositions~3.2 and~3.3
state, without proofs, that each step preserves and reflects satisfiability
and that a satisfiability-equivalent simple system is obtained in linear time
\cite[\S3, p.~10]{SchmidtSchaussSmolka1991}.

The source then associates a graph, its \emph{skeleton}, with a simple system:
concept variables are nodes; existential and universal constraints are
directed edges; a union constraint joins its source to its two alternatives;
and atomic or complemented-atomic inclusions label nodes.  Individual
constraints are not part of this graph.  A constraint tree has a tree
skeleton.  It is \emph{fresh} exactly when it results from unfolding one
simple description; it has one individual constraint \(x:X\) at the skeleton
root.  Theorem~3.4 composes normalization, Proposition~3.1, and unfolding:
from a description \(C\), a fresh tree satisfiable exactly when \(C\) is
coherent is computable in linear time
\cite[pp.~10--11]{SchmidtSchaussSmolka1991}.

\subsubsection{Clashes, completion, and the extracted interpretation}

For a simple system, the article defines a \emph{clash} as either
\[
 \{x:X,\ X\sqsubseteq\neg\top\}\subseteq S
\]
or, for a concept symbol \(A\),
\[
 \{x:X,\ X\sqsubseteq A,\ x:Y,\ Y\sqsubseteq\neg A\}\subseteq S.
\]
This is a condition on generated constraints.  It is not merely an informal
claim that two arbitrary complex formulas look contradictory.
Proposition~4.1 states, without proof, that every system containing a clash is
unsatisfiable and that a clash can be found in linear time
\cite[\S4, p.~11]{SchmidtSchaussSmolka1991}.

A simple system \(S\) is \emph{complete} exactly when all three requirements
hold:
\begin{enumerate}
\item if \(x:X\) and \(X(\exists R)Y\) belong to \(S\), some \(y\) has both
      \(xRy\) and \(y:Y\) in \(S\);
\item if \(x:X\), \(X(\forall R)Y\), and \(xRy\) belong to \(S\), then
      \(y:Y\) belongs to \(S\);
\item if \(x:X\) and \(X\sqsubseteq Y\sqcup Z\) belong to \(S\), then
      \(x:Y\) or \(x:Z\) belongs to \(S\).
\end{enumerate}
Proposition~4.2 says that a complete system is satisfiable exactly when it is
clash-free.  For the nontrivial direction, let \(V\) be the finite set of
individual variables occurring in \(S\), and let \(\ast\) be a fresh element.
Define
\[
\begin{aligned}
 D^{\mathcal I}
   &=
   \begin{cases}
    V,&V\ne\varnothing,\\
    \{\ast\},&V=\varnothing,
   \end{cases}\\
 A^{\mathcal I}
   &=\{x\in V:\exists X\,
          (x:X\in S\land X\sqsubseteq A\in S)\}
      &&(A\ne\top),\\
 R^{\mathcal I}
   &=\{(x,y):xRy\in S\},\\
 \top^{\mathcal I}&=D^{\mathcal I},\qquad
 \alpha(X)=\{x\in V:x:X\in S\}.
\end{aligned}
\]
Choose \(x_0\in D^{\mathcal I}\), and define the assignment on every
individual variable by
\[
 \alpha(x)=
 \begin{cases}
  x,&x\in V,\\
  x_0,&x\notin V.
 \end{cases}
\]
The domain is nonempty and the individual assignment is total.
No constraint mentions an individual variable outside \(V\).
If \(V=\varnothing\), every ordinary concept symbol, role, and concept
variable has empty extension, and every constraint is satisfied vacuously.
A fresh tree has an individual root, so its \(V\) is nonempty.

For general \(V\), individual memberships and role edges hold by
definition.  An atomic inclusion holds by the definition of
\(A^{\mathcal I}\).  A complemented-atomic inclusion can fail only if the
second clash pattern is present.  If \(x:X\in S\) and
\(X\sqsubseteq\neg\top\in S\), the first clash pattern is present.  The three
remaining constraint forms hold by completeness conditions 1--3,
respectively.  Thus \((\mathcal I,\alpha)\) satisfies \(S\).
Conversely, either clash pattern forces an element to satisfy both a
concept and its complement, which is impossible.
This gives the canonical-model argument for
Proposition~4.2, including the total assignment and every constraint form
\cite[Prop.~4.2, p.~11]{SchmidtSchaussSmolka1991}.

Before imposing control, Proposition~4.3 states the exact semantic effects of
the three kinds of completion:
\begin{enumerate}
\item if \(x:X,X(\exists R)Y\in S\) and \(y\) is fresh, then \(S\) is
      satisfiable exactly when
      \(S\cup\{xRy,y:Y\}\) is satisfiable;
\item if \(x:X,xRy,X(\forall R)Y\in S\), then \(S\) is satisfiable exactly
      when \(S\cup\{y:Y\}\) is satisfiable;
\item if \(x:X,X\sqsubseteq Y\sqcup Z\in S\), then \(S\) is satisfiable
      exactly when at least one of
      \(S\cup\{x:Y\}\) and \(S\cup\{x:Z\}\) is satisfiable.
\end{enumerate}
The article prints no proof of the proposition.  For item~1 a satisfying
assignment is extended by an existential witness; item~2 is the universal
constraint clause; and item~3 is membership in a union.  In every reverse
direction, removing newly added constraints preserves satisfaction.  The
third item is a disjunction between two possible systems: it does not say that
an arbitrary chosen branch preserves satisfiability
\cite[Prop.~4.3, p.~12]{SchmidtSchaussSmolka1991}.

Figure~1 turns those facts into controlled rules.  The existential rule adds
\(xRy,y:Y\), with \(y\) fresh, only when
\[
 x:X,\ X(\exists R)Y\in S
 \quad\text{and no }z\text{ has both }xRz,z:Y\in S.
\]
The universal rule adds \(y:Y\) only when
\[
 x:X,\ xRy,\ X(\forall R)Y\in S
 \quad\text{and }y:Y\notin S.
\]
The union rule, when \(x:X\) and \(X\sqsubseteq Y_1\sqcup Y_2\) are present
but neither alternative membership is present, chooses one
\(Z\in\{Y_1,Y_2\}\) and adds \(x:Z\).  Only this last rule branches.
Proposition~4.4 states, again without proof, that no infinite basic-rule chain
starts at a fresh tree, that normal forms are exactly complete systems, and
that one basic step from a constraint tree produces another constraint tree
and satisfies the backward implication
\[
 T'\text{ satisfiable}\Longrightarrow T\text{ satisfiable}
 \qquad(T\to T').
\]
The last implication is the direction valid for every chosen union branch
\cite[Fig.~1 and Prop.~4.4, p.~12]{SchmidtSchaussSmolka1991}.

A \emph{completion} of a constraint tree \(T\) is a complete constraint tree
obtainable from \(T\) by finitely many applications of the basic completion
rules.  Theorem~4.5 states
\[
 \boxed{
 \text{a fresh tree is satisfiable}
 \quad\Longleftrightarrow\quad
 \text{it has a clash-free completion}.}
\]
It also states that every fresh tree has a completion.  The theorem's entire
printed proof is one sentence referring to the preceding propositions.  The
expanded dependency argument is as follows.  Termination and the normal-form
criterion give a completion.  If a completion is clash-free,
Proposition~4.2 makes it satisfiable and Proposition~4.4 propagates
satisfiability backward.  Conversely, from a satisfiable tree,
Proposition~4.3 permits a satisfiable union alternative whenever one is
needed; termination yields a satisfiable completion, which
Proposition~4.1 makes clash-free.  This is a handbook reconstruction of the
one-sentence proof boundary, not a claim that the journal printed these
steps
\cite[Thm.~4.5, p.~13]{SchmidtSchaussSmolka1991}.

Two handbook-supplied rule traces make the mechanism explicit.  The
description
\[
 \exists r:A\sqcap\forall r:\neg A
\]
forces an \(r\)-successor through the existential edge and then propagates
\(\neg A\) to that same successor through the universal edge; its \(A\) and
\(\neg A\) labels clash.  By contrast,
\[
 A\sqcap\exists r:B\sqcap\forall r:B
\]
has a clash-free completion with a root in \(A\), one \(r\)-successor in
\(B\), and no other positive atomic requirements.  These examples apply the
source rules but are not examples printed in the article.  Its Figure~2
instead exhibits a recursively growing, one-role family for which every
completion has exponentially many individual variables; the caption gives a
pattern with an ellipsis rather than a formally indexed family
\cite[Fig.~2, p.~13]{SchmidtSchaussSmolka1991}.

\subsubsection{Traces and the linear-space upper bound}

A \emph{partial completion} of a fresh tree is any tree reachable by finitely
many basic steps.  Given a constraint tree \(T\) and individual variables
\(x,y\) occurring in \(T\), \(y\) is a \emph{successor} of \(x\), and \(x\) a
\emph{predecessor} of \(y\), exactly when \(xRy\in T\) for some role \(R\).
An occurring individual variable \(x\) is an \emph{individual root} exactly
when \(x:X\in T\) and \(X\) is the root of the skeleton of \(T\).  These are
relations and a distinguished variable among the individuals represented by
the constraints, not the skeleton's edges among concept variables;
Proposition~5.1 needs this distinction.  It states that every partial
completion has one individual root,
that the root has no predecessor, and that every other individual variable has
exactly one predecessor.  A partial completion \(U\) is a \emph{trace} exactly
when
\begin{enumerate}
\item no individual variable in \(U\) has more than one successor in total;
\item neither the universal rule nor the union rule applies; and
\item every applicable basic existential step would give some individual a
      second successor.
\end{enumerate}
The restricted trace existential rule creates \(xRy,y:Y\) only when
\[
 x:X,\ X(\exists R)Y\in S
 \quad\text{and there are no }R',y'\text{ with }xR'y'\in S.
\]
The variable \(y\) must not occur in \(S\).  Its no-successor condition ranges
over every role, not merely over \(R\).
Together with the unchanged universal and union rules it forms the source
relation \(\to_T\).  Traces are exactly the reachable \(\to_T\)-normal forms
\cite[\S5, pp.~14--15]{SchmidtSchaussSmolka1991}.

Proposition~5.2 states four facts and supplies no printed proofs:
\begin{enumerate}
\item a \(\to_T^*\)-derivative never associates two distinct individual
      variables with the same concept variable;
\item every such derivation has length linear in the initial tree;
\item every trace is contained in a completion;
\item every completion is the union of finitely many traces.
\end{enumerate}
Figure~3 assumes that individual variables are natural numbers and defines
\[
 \operatorname{eval}:
 \mathsf{nat}\times\{\text{constraint trees}\}\longrightarrow\mathsf{bool}.
\]
At the current individual \(x\), it first returns false if either clash pattern
is present at \(x\).  If
\[
 x:X,\quad X\sqsubseteq Y\sqcup Z\in S,
 \qquad x:Y\notin S,\quad x:Z\notin S,
\]
it returns the disjunction of the calls on
\(S\cup\{x:Y\}\) and \(S\cup\{x:Z\}\).  Otherwise it puts \(y=x+1\).  For
each \(x:X,X(\exists R)Y\in S\), it recursively evaluates
\[
 \operatorname{eval}\!\left(
  y,
  S\cup\{y:Y\}\cup
  \{y:Z'\mid \exists Z\,(x:Z\in S\land Z(\forall R)Z'\in S)\}
 \right),
\]
where the role \(R\) in the universal constraint is the role of that selected
existential edge, and it conjoins the answers for all such edges.  The empty
conjunction is true.  The displayed source call does not add \(xRy\): it moves
the current individual to \(y\) and retains only the successor memberships
needed for that trace.  It is therefore not a literal basic completion step,
whose existential rule does add \(xRy\).  For the individual root \(x\) of a
fresh tree \(T\), the figure asserts
\[
 \operatorname{eval}(x,T)=\mathsf{true}
 \quad\Longleftrightarrow\quad
 T\text{ has a clash-free completion}.
\]

The recursion depth is at most the height of the original constraint tree.
Apart from control information, only one trace need be retained at once;
Proposition~5.2 bounds that trace linearly in the tree, and Theorem~3.4 bounds
the tree linearly in the input description.  Theorem~5.3 therefore gives a
linear-space decision procedure for \(\mathcal{ALC}\) coherence and
subsumption.  The article describes the procedure as exponential-time but
does not print a machine encoding, bit-level stack analysis, or explicit
constant.  Its statement that total correctness follows from the preceding
propositions is not a substitute for a separately printed induction
\cite[Fig.~3 and Thm.~5.3, p.~15]{SchmidtSchaussSmolka1991}.

The source uses a visually similar but different restricted existential rule
for the smaller languages:
\[
 S\to_{\top\exists}S\cup\{xRy,y:Y\}
\]
when \(x:X,X(\exists R)Y\in S\), no individual \(y'\) has \(xRy'\in S\),
and \(y\) is fresh.  The side condition fixes the same role \(R\), so the rule
may create different successors for different role symbols.  The trace rule
\(\to_{T\exists}\), by contrast, requires that \(x\) have no successor by any
role.  Hence
\[
 \to_{T\exists}\subseteq\to_{\top\exists},\qquad
 \to_{\top}
   =\to_{\top\exists}\cup\to_{\forall}\cup\to_{\sqcup},
 \qquad
 \to_T\subseteq\to_{\top}.
\]
The \(\top\)-completions of a fresh constraint tree \(T\) are the reachable
\(\to_{\top}\)-normal forms of \(T\).  The printed subscripts \(T\) and
\(\top\) must not be conflated.

The restricted-language bounds use four intermediate results.  Proposition
5.4 says that, for a fresh tree without union constraints, all completions
agree up to consistent renaming and the following are equivalent:
\[
\begin{gathered}
 T\text{ is satisfiable};\qquad
 T\text{ has a clash-free completion};\\
 \text{every completion of }T\text{ is clash-free};\qquad
 \text{every trace of }T\text{ is clash-free}.
\end{gathered}
\]
Proposition~5.6 says that every \(\to_{\top}\)-derivation from a fresh tree
has length linear in that tree.  Its printed proof gives the invariant that at
most one individual variable is associated with each concept variable and
then calls the conclusion obvious.

Let \(U\) be a \(\top\)-completion of a fresh constraint tree \(T\), where
\(T\) is obtainable by unfolding an \(\mathcal{ALU}\) description.
Proposition~5.7 states that
\[
\begin{aligned}
 &x:X,\ X(\exists R)Y,\ xRy\in U\\
 &\qquad\Longrightarrow\quad
 \bigl(U\text{ satisfiable}\Longleftrightarrow
       U\cup\{y:Y\}\text{ satisfiable}\bigr).
\end{aligned}
\]
If \(y:Y\in U\), the two systems displayed there are identical; the
proposition's nontrivial, missing-target instance therefore has the additional
condition \(y:Y\notin U\).  Its one-sentence proof uses the fact that an
existentially reachable concept variable is an unlabeled leaf.
Proposition~5.8 extends a clash-free
\(\top\)-completion to a basic completion through Proposition~5.7 and obtains
\[
 T\text{ is satisfiable}
 \quad\Longleftrightarrow\quad
 T\text{ has a clash-free }\top\text{-completion}
\]
for every fresh tree \(T\) unfolded from an \(\mathcal{ALU}\) description.

The ensuing exact bounds are:
\begin{itemize}
\item \(\mathcal{ALE}\)-incoherence is decidable in nondeterministic linear
      time, so \(\mathcal{ALE}\)-coherence is in co-NP
      (Theorem~5.5);
\item \(\mathcal{ALU}\)-coherence is NP-complete
      (Theorem~5.9);
\item \(\mathcal{AL}\)-coherence is decidable in linear time
      (Theorem~5.10).
\end{itemize}
The middle result combines Proposition~2.3, Proposition~5.6, and
Proposition~5.8: guess a linear \(\top\)-completion and check it for a clash.
For the last result, all \(\top\)-completions agree up to consistent
individual-variable renaming; therefore any one completion and one clash test
suffice
\cite[pp.~15--17]{SchmidtSchaussSmolka1991}.

Nothing in this termination or space argument is blocking.  The source uses
completion rules, constraint trees, and traces; it defines no tableau node,
branch label, ancestor-blocking relation, anywhere-blocking relation, caching
condition, or cyclic-TBox procedure.  Space is controlled by discarding
completed traces, not by identifying a generated node with an ancestor.

\subsubsection{The QBF lower bound}

Section~6 reduces a source-defined quantified Boolean validity problem.  A
literal is a nonzero integer; a clause is a nonempty sequence
\(l_1\ldots l_k\) ordered by nondecreasing absolute value; and, for positive
integers \(m\le n\), a prefix from \(m\) through \(n\) is
\[
 (Q_m m)(Q_{m+1}(m+1))\cdots(Q_n n),
 \qquad Q_i\in\{\exists,\forall\}.
\]
An input \(P.M\) has a prefix \(P\) from \(1\) through \(n\), a finite
nonempty set \(M\) of clauses, and only literals
\[
 \{-n,\ldots,-1,1,\ldots,n\}
\]
in those clauses.  A \(P\)-assignment \(\sigma\) maps the variables of \(P\)
to \(\{0,1\}\), with signed-literal and clause satisfaction defined by
\[
\begin{aligned}
 \sigma\models l
 &\iff
   (l>0\land\sigma(l)=1)
   \ \lor\
   (l<0\land\sigma(-l)=0),\\
 \sigma\models l_1\ldots l_k
 &\iff
   \exists i\in\{1,\ldots,k\}\;(\sigma\models l_i).
\end{aligned}
\]
A nonempty
set \(\mathcal A\) of \(P\)-assignments is \emph{canonical for \(P\)} by the
following recursion:
\begin{itemize}
\item if \(P=(\exists m)P'\), all assignments agree at \(m\), and, when the
      tail \(P'\) is nonempty, their tail restrictions form a canonical set
      for \(P'\);
\item if \(P=(\forall m)P'\), both values at \(m\) occur, and, when \(P'\) is
      nonempty, the restrictions of each value class form a canonical set for
      \(P'\).
\end{itemize}
Thus the term names a set-of-assignments representation of the quantifier
choices, not a new object-language connective.  A formula \(P.M\) is valid
exactly when some canonical set for \(P\) consists entirely of assignments
satisfying every clause in the finite nonempty set \(M\).  Theorem~6.1 cites
the PSPACE-completeness of this problem; the article does not prove it
\cite[\S6.1, p.~18]{SchmidtSchaussSmolka1991}.

Fix one role symbol \(R\) and one concept symbol \(A\).  For
\(M=\{C_1,\ldots,C_k\}\), put
\[
 [P.M]=[P]\sqcap[C_1]^0\sqcap\cdots\sqcap[C_k]^0.
\]
The prefix translation is
\[
\begin{aligned}
 [(\exists m)P]
   &=((\exists R:A)\sqcup(\exists R:\neg A))
      \sqcap(\forall R:[P]),\\
 [(\forall m)P]
   &=(\exists R:A)\sqcap(\exists R:\neg A)
      \sqcap(\forall R:[P]),\\
 [(\exists m)]
   &=(\exists R:A)\sqcup(\exists R:\neg A),\\
 [(\forall m)]
   &=(\exists R:A)\sqcap(\exists R:\neg A).
\end{aligned}
\]
An existential quantifier selects one truth-value successor through a union.
A universal quantifier requires both truth-value successors.  In either
nonterminal case, the universal restriction sends the remaining prefix to
every successor.

For an ordered clause, write \(lC\) for first literal \(l\) followed by the
remaining clause \(C\).  The level-indexed clause translation is
\[
\begin{aligned}
 [lC]^m&=\forall R:[lC]^{m+1}&&(|l|>m),\\
 [mC]^m&=A\sqcup[C]^m,&
 [-mC]^m&=\neg A\sqcup[C]^m,\\
 [l]^m&=\forall R:[l]^{m+1}&&(|l|>m),\\
 [m]^m&=A,&[-m]^m&=\neg A.
\end{aligned}
\]
Unfolding \([P]\) produces the \emph{prefix tree}; unfolding
\([C_i]^0\) produces the \(i\)-th \emph{clause tree}.  These names distinguish
the source-defined components of the one fresh constraint tree for
\([P.M]\).  A \emph{prefix completion} is a completion of the prefix tree.
The index delays a literal until the level of its variable.  Every union
introduced in a clause tree consequently has the restricted shape
\(A\sqcup D\) or \(\neg A\sqcup D\).  The whole translation is computable in
quadratic time
\cite[\S6.2, pp.~19--21]{SchmidtSchaussSmolka1991}.

The article assigns the skeleton root level \(0\); an existential or universal
edge raises the level by one; all variables in one union constraint share a
level; \(x:X\) gives \(x\) the level of \(X\); and \(xRy\) puts \(y\) one level
above \(x\).  These clauses uniquely level every tree reachable from a fresh
tree by completion steps.  A trace through the prefix tree has one individual
at each level.  Membership through a node labelled \(A\) records value \(1\),
and membership through a node labelled \(\neg A\) records value \(0\).

A clause tree has no existential edge.  Its leaves correspond one-for-one to
the literals of its clause: positive \(l\) gives an \(A\)-labelled leaf at
level \(l\), while negative \(l\) gives a \(\neg A\)-labelled leaf at level
\(-l\).  Completing the tree's union chain selects one such leaf along a
trace.  This level-and-label correspondence is what lets a clash test express
literal satisfaction.

Lemma~6.2 states the reduction invariant
\[
 \boxed{
 P.M\text{ is valid}
 \quad\Longleftrightarrow\quad
 [P.M]\text{ is coherent}.}
\]
Every prefix completion is clash-free.  Its traces represent a canonical set
of assignments: existential union keeps one value at its level, universal
intersection creates both values, and the trailing universal restriction
recurses through the remaining prefix.  Conversely, the same equations turn
each canonical set into a prefix completion.  The article identifies this
correspondence as crucial but does not print the quantifier-by-quantifier
induction; that induction is the first reconstructed step in the account
below.

For coherence implying validity, take a clash-free completion of the whole
translated tree and its prefix completion.  Fix a represented prefix trace.
For each clause, a trace through the attached clause tree selects one literal
leaf.  Clash-freedom forces that leaf's \(A\) or \(\neg A\) label to agree with
the prefix value at the same level.  Hence the represented assignment satisfies
that clause.  Every assignment represented by the prefix completion satisfies
every clause, so the represented canonical set witnesses validity.

For validity implying coherence, start with a canonical satisfying set and the
corresponding prefix completion.  The source restricts the union rule during
clause completion as follows.  When
\[
 X\sqsubseteq Y_1\sqcup Y_2,\quad
 x:X,\quad x:Z,\quad Z\sqsubseteq C,\quad Y_1\sqsubseteq D
\]
are present, with \(x:Y_1\notin S\) and \(x:Y_2\notin S\), and where
\(C,D\in\{A,\neg A\}\) are the already represented literal labels, add
\(x:Y_1\) if \(C=D\) and \(x:Y_2\) otherwise.  This relation, denoted
\(\to_Q\), is a deterministic restriction of the basic union rule: it chooses
the displayed literal when it agrees with the fixed prefix value and continues
along the clause when it does not.  Universal propagation and \(\to_Q\)
complete each clause tree without creating existential edges.
If a clash first appeared at the greatest selected literal level, a prefix
trace through the same individual chain would represent an assignment
satisfying no literal of that clause, contrary to the chosen canonical set.
The source prints both directions of this argument but calls the canonical-set
correspondence crucial without proving its induction and calls the final
literal check easy.  The preceding expansion marks those two proof obligations
rather than attributing their details to the journal
\cite[Lem.~6.2, pp.~20--21]{SchmidtSchaussSmolka1991}.

The article's concrete example is
\[
 \forall x\,\exists y\,
 ((\neg x\lor y)\land(x\lor\neg y)).
\]
It is valid by choosing \(y=x\).  In the source's Figure~4 translation, the
universal variable creates both an \(A\)-successor and a \(\neg A\)-successor
at level \(1\).  At level \(2\), the existential union chooses the matching
truth value for \(y\) on each branch.  The two clause trees are aligned to
those levels, so each prefix trace can select a leaf whose label agrees with
the corresponding truth value.  This is an independently worded account of
the source example; the journal figure itself is not reproduced
\cite[Fig.~4, pp.~18--19]{SchmidtSchaussSmolka1991}.

Theorem~6.3 combines Lemma~6.2, the quadratic translation, cited QBF hardness,
and Theorem~5.3:
\[
 \boxed{\text{\(\mathcal{ALC}\) concept coherence and subsumption are
 PSPACE-complete.}}
\]
Coherence hardness is direct.  The printed subsumption-hardness sentence
compresses one closure step.  For any \(L\in\mathsf{PSPACE}\), closure under
complement gives \(\overline L\in\mathsf{PSPACE}\).  Since QBF validity is
PSPACE-complete, there is a polynomial reduction \(f\) with
\[
 x\in\overline L
 \quad\Longleftrightarrow\quad
 f(x)\text{ is QBF-valid}.
\]
Therefore \(x\in L\) exactly when \(f(x)\) is QBF-invalid, so QBF invalidity
is PSPACE-hard.  Lemma~6.2 maps invalidity to incoherence, and
Proposition~2.1 maps incoherence linearly to the subsumption query
\(C\sqsubseteq\neg\top\).  This complement-reduction step completes the
printed argument but is not written out in the article
\cite[Thm.~6.3, p.~21]{SchmidtSchaussSmolka1991}.

\subsubsection{Source boundary and later knowledge-base calculi}

The proof granularity matters.  Proposition~2.3 prints both directions of its
CNF reduction, and Theorem~3.4 prints a short composition of normalization,
translation, and unfolding.  By contrast, Propositions~2.1--2.2, 3.1--3.3,
4.1, 4.3--4.4, 5.1--5.2, and~5.4 are stated without proofs.
Proposition~4.2's canonical model and total assignment are given above,
with verification of every constraint form.  Theorem~4.5 has
only its one-sentence dependency reference.  The article's exact Theorem~4.6
states
\[
 \text{Every \(\mathcal{FL}\) concept description is coherent.}
\]
Here \(\mathcal{FL}\) is retained only inside the source theorem as the
article's identifier for an earlier formal-language presentation.  The article
refers that language's full grammar to prior sources rather than defining it in
full, so the owed artifact is a separate source presentation of those
authorities; the handbook does not introduce \(\mathcal{FL}\) as a usable
language here, and no later result depends on it.  The short proof translates
into \(\mathcal{ALC}\) and always selects right-hand union alternatives
\cite[Thm.~4.6, p.~13]{SchmidtSchaussSmolka1991}.

Theorem~5.3 is preceded by the trace-space explanation but has no machine-level
accounting.  Theorem~5.5 has a short nondeterministic-trace proof.
Proposition~5.6 prints one invariant and calls the conclusion obvious;
Proposition~5.7 has a one-sentence unlabeled-leaf proof; and
Proposition~5.8 has a short completion-extension proof.  Theorems~5.9 and~5.10
also have short proofs.  Theorem~6.1 is cited rather than proved.
Lemma~6.2 prints both directions while leaving the canonical correspondence
and final literal induction abbreviated, and Theorem~6.3 omits the explicit
complement reduction.  The reconstructions above fill only those identified
gaps.

The constraint grammar distinguishes individual membership \(x:X\) from
concept inclusion \(X\sqsubseteq A\) or \(X\sqsubseteq\neg A\).
The canonical construction above interprets these forms with a total
individual assignment.  The trace subscript \(T\) and the top symbol
\(\top\) also denote rules with different successor conditions, as
specified in their respective rule displays.

Section~7 separately compares general roles with functional roles, called
features.  It notes distributivities valid for functions but invalid for
arbitrary binary relations and cites neighboring feature-description results.
It does not extend the concept-coherence calculus to a TBox, ABox, role
hierarchy, number restriction, or blocking procedure
\cite[\S7, pp.~22--24]{SchmidtSchaussSmolka1991}.

\paragraph{The two registered book presentations.}
The 2003 first edition of \emph{The Description Logic Handbook} is an
independently registered source.  Only its front matter and later-edition
chapter metadata are available in the present source basis, so its exact
knowledge-base syntax, rule tables, blocking condition, complexity theorems,
and proofs are not stated here or reconstructed from the 2007 revision or the
2017 textbook \cite{BaaderEtAl2003}.

The authoritative 2017 chapter record, without the chapter body, supports only
this summary-level classification: concept satisfiability and subsumption
without a supplied TBox are PSPACE-complete, whereas the same tasks relative
to an arbitrary finite TBox as in
Definition~\ref{def:p09-alc} are EXPTIME-complete.  The record also says that
the bounds transfer to satisfiability of the corresponding finite knowledge
bases \cite{BaaderHorrocksLutzSattler2017}.  The exact blocking rules,
termination measure, worked example, and proof organization are not established
by that record and are not supplied by the earlier concept calculus.

The blocking procedure in Definition~\ref{def:p09-alc-tableau} has its
own correctness and termination proof in Annex~\ref{annex:p9.alc.1}.
Section~\ref{sec:p09-alc-general-tboxes} additionally supplies complete
concept types, elimination of unrealizable types, a finite canonical model,
and a combined-complexity analysis.  Theorem~\ref{thm:p09-alc-general-exptime}
and Annex~\ref{annex:p9.alc.3} give the general-TBox classification with
an explicit alternating-machine lower bound and reductions for the finite
knowledge-base tasks.

The acquired sources for that classification are Donini and Massacci's
2000 paper and the author draft of Baader, Horrocks, and Sattler's
2008 chapter.  The former has its own tableau search techniques and a
concept/TBox input; the latter uses looping tree automata
\cite[Theorem~3.13, p.~98]{DoniniMassacci2000}
\cite[draft Theorem~6, p.~26]{BaaderHorrocksSattler2008}.
Those sources and the handbook's explicit constructions have their own
attribution.  The exact 2003 first-edition and 2017 textbook presentations
still require their respective bodies, including their rule tables,
blocking conditions, examples, and proof organization.
